\documentclass[11pt,a4paper]{article}

\usepackage[main=italian,provide=*]{babel}
\usepackage{amsmath,amssymb,mathtools}
\usepackage{booktabs}
\usepackage{colortbl}
\usepackage{graphicx}
\usepackage{geometry}
\geometry{margin=2.6cm}
\usepackage{xcolor}
\usepackage{fancyhdr}
\usepackage{enumitem}

\newcommand{\ket}[1]{\lvert #1\rangle}
\newcommand{\vs}[1]{\mathbf{s}_{#1}}
\newcommand{\vS}{\mathbf{S}}

\definecolor{colA}{RGB}{230,126,34}
\definecolor{colB}{RGB}{41,128,185}
\definecolor{colC}{RGB}{39,174,96}
\definecolor{optA}{RGB}{39,174,96}
\definecolor{optB}{RGB}{192,57,43}

\pagestyle{fancy}
\fancyhf{}
\fancyhead[L]{\small Trimero isoscele — analisi del termine DM}
\fancyhead[R]{\small\thepage}
\renewcommand{\headrulewidth}{0.3pt}

\title{\bfseries Il termine DM sul trimero isoscele:\\
perché una scelta apre il gap all'incrocio e l'altra no}
\author{Note di studio — Tesi triennale, Samuele Schiavi\\ \small Relatore: Prof.\ Alessandro Chiesa}
\date{}

\begin{document}
\maketitle
\thispagestyle{fancy}

\begin{abstract}
\noindent
Questo documento ricostruisce per intero l'analisi che ha portato a una conclusione
non ovvia: fra due forme plausibili del termine di Dzyaloshinskii--Moriya (DM) sul
trimero isoscele, \emph{una sola} apre effettivamente il gap all'incrocio di livello
$b_c$, mentre l'altra — pur sembrando ragionevole e pur avendo un residuo numerico
non nullo se calcolato nel punto sbagliato — lascia l'incrocio \emph{esatto}. La
spiegazione finale è una regola di simmetria precisa (non un dettaglio numerico), e il
documento la deriva passo per passo, includendo l'errore metodologico in cui siamo
incappati durante l'analisi e come lo abbiamo scoperto.
\end{abstract}

\section{Il punto di partenza: cosa già sappiamo}

Nella teoria esatta del trimero isoscele (decomposizione di Kambe), l'Hamiltoniana
\[
H_0 = J\,\boldsymbol\sigma_1\cdot\boldsymbol\sigma_2
+ J'(\boldsymbol\sigma_2\cdot\boldsymbol\sigma_3+\boldsymbol\sigma_3\cdot\boldsymbol\sigma_1)
+ b\sum_{i=1}^3 Z_i
\]
si risolve in forma chiusa grazie alla conservazione di $\mathbf S_{12}^2$ (simmetria
di scambio $1\leftrightarrow2$) oltre a $\mathbf S^2,S_z$. Tre multipletti: $A$
($S_{12}{=}1,S{=}\tfrac32$), $B$ ($S_{12}{=}1,S{=}\tfrac12$), $C$
($S_{12}{=}0,S{=}\tfrac12$). Al punto di lavoro scelto ($J{=}1,J'{=}0.4$), il
fondamentale passa da $C$ ad $A$ a $b_c=2J+J'=2.4$: un \textbf{incrocio vero}, non un
anticrossing — le due curve di energia $E_C(b)=-3J-b$ e $E_A(b)=J+2J'-3b$ si
attraversano senza repulsione, verificato a precisione macchina.

\textbf{La domanda di questo documento.} Se si aggiunge un termine DM (che in generale
rompe $M$, e quindi collega stati con $\Delta M=\pm1$), questo incrocio si trasforma
in un anticrossing — cioè il gap si apre — per \emph{qualunque} forma di DM, o solo
per alcune?

\section{La forma generale del DM e la sua riduzione}

L'interazione DM fra due spin è $\vec D_{ij}\cdot(\vs i\times\vs j)$, con $\vec
D_{ij}$ un vettore (pseudo-vettore) che dipende dalla geometria del legame. Nel lavoro
sul dimero, per una singola componente non nulla, questo si riduce a
\[
H_{DM}^{(ij)} = D\,(X_iZ_j - Z_iX_j),
\]
la stessa forma reale già validata e discussa in \texttt{simmetrie\_dmi\_spin.pdf}. Per
il trimero, la generalizzazione naturale è un termine di questo tipo su \emph{ciascun}
legame:
\[
H_{DM} = D_{12}(X_1Z_2-Z_1X_2) + D_{23}(X_2Z_3-Z_2X_3) + D_{31}(X_3Z_1-Z_3X_1).
\]
Restano da fissare i tre coefficienti $D_{12},D_{23},D_{31}$.

\section{Vincolo di simmetria: quali disposizioni sono ammissibili}

\subsection{Perché serve un vincolo}

L'intera teoria del trimero isoscele assume che l'Hamiltoniana sia invariante per lo
scambio dei siti $1,2$ (rappresentato dall'operatore di permutazione $P_{12}$ sui
qubit). Se aggiungiamo $H_{DM}$ senza vincoli, rischiamo di introdurre una rottura di
questa simmetria non giustificata dalla geometria del problema — un'incoerenza interna
rispetto a tutto il resto della teoria costruita sopra quell'ipotesi.

\subsection{Verifica esplicita, caso per caso}

Ho costruito esplicitamente $P_{12}$ come matrice di permutazione $8\times8$ e
verificato $\|P_{12}H_{DM}P_{12}^\dagger - H_{DM}\|$ per tutte le combinazioni di
segno di $(D_{12},D_{23},D_{31})\propto(J,J',J')$, a $r=0.15$:

\begin{center}
\renewcommand{\arraystretch}{1.25}
\begin{tabular}{cccc}
\toprule
$\mathrm{segno}(D_{12})$ & $\mathrm{segno}(D_{23})$ & $\mathrm{segno}(D_{31})$ & $\|[P_{12},H_{DM}]\|$ \\
\midrule
$0$ & $+$ & $-$ & $\mathbf{0}$ \quad (simmetrica) \\
$0$ & $-$ & $+$ & $\mathbf{0}$ \quad (simmetrica) \\
$0$ & $+$ & $+$ & $0.24$ \\
$+$ & $+$ & $+$ & $0.42$ \\
$+$ & $+$ & $-$ & $0.30$ \\
\ldots & & & (tutte le altre combinazioni: $\neq0$) \\
\bottomrule
\end{tabular}
\end{center}

\textbf{Risultato netto:} fra tutte le disposizioni testate, \emph{una sola} rispetta
esattamente la simmetria di scambio: \textbf{nessun DM sul legame di base}
($D_{12}=0$), con segno \textbf{opposto} sui due legami laterali equivalenti
($D_{23}=-D_{31}\equiv D'$).

\subsection{Le due opzioni che restano sul tavolo}

\begin{center}
\renewcommand{\arraystretch}{1.4}
\begin{tabular}{@{}p{2.3cm}p{6.3cm}p{6.3cm}@{}}
\toprule
& \textbf{Opzione A} (simmetrica) & \textbf{Opzione B} (proposta del relatore) \\
\midrule
forma & $D'\big[(X_2Z_3{-}Z_2X_3)-(X_3Z_1{-}Z_3X_1)\big]$ &
$r\big[J(X_1Z_2{-}Z_1X_2)+J'(X_2Z_3{-}Z_2X_3)+J'(X_3Z_1{-}Z_3X_1)\big]$ \\
simmetria $P_{12}$ & rispettata esattamente & rotta \\
$[\,H,\mathbf S_{12}^2\,]$ & $=0$ esatto (verificato, $\|\cdot\|\sim10^{-16}$) & $\neq0$
($\|\cdot\|\sim1.4$ a $r{=}0.15$) \\
struttura residua & blocco $2{\times}2$ ($C$) $+$ blocco $6{\times}6$ ($A{+}B$ misti) &
nessuna, $8\times8$ pieno \\
\bottomrule
\end{tabular}
\end{center}

\section{Un errore metodologico, e come lo abbiamo scoperto}

\subsection{Il controllo sbagliato: gap a $b_c$ fisso}

Il primo controllo fatto è stato calcolare il gap fra i due autovalori più bassi
\emph{esattamente} al vecchio punto critico $b_c=2.4$ (quello di $D=0$), per diversi
valori del parametro DM:

\begin{center}
\renewcommand{\arraystretch}{1.3}
\begin{tabular}{cccc}
\toprule
parametro & gap Opzione A a $b{=}2.4$ & gap Opzione B a $b{=}2.4$ \\
\midrule
$0.05$ & $0.0058$ & $0.0853$ \\
$0.15$ & $0.0579$ & $0.2612$ \\
$0.30$ & $0.1702$ & $0.4577$ \\
\bottomrule
\end{tabular}
\end{center}

A prima vista, sembra che \emph{entrambe} le opzioni aprano il gap — solo l'Opzione B
lo apre di più. \textbf{Questa conclusione è sbagliata.}

\subsection{Cosa c'è di sbagliato}

Calcolare il gap a $b=2.4$ presuppone implicitamente che il punto di incrocio resti
esattamente lì anche con il DM acceso. Non c'è alcuna ragione per cui debba essere
così: aggiungere un termine all'Hamiltoniana può spostare il punto in cui le energie
si toccano, non solo aprirle o lasciarle chiuse. Un valore di gap positivo calcolato
in un punto che non è più il vero punto critico non dice nulla sull'apertura
dell'incrocio — dice solo che le due curve, spostato il punto di confronto, non sono
più esattamente sovrapposte lì.

\subsection{Il controllo corretto: minimizzare il gap su tutto $b$}

Il controllo corretto è cercare, per ciascun valore del parametro DM, il \emph{minimo
assoluto} del gap su tutto l'asse $b$ (non il valore a $b=2.4$ fisso), usando
un'ottimizzazione numerica di precisione:
\[
g_\text{min}(D) = \min_b\ \big[E_1(b;D) - E_0(b;D)\big].
\]

\begin{center}
\renewcommand{\arraystretch}{1.3}
\begin{tabular}{ccccc}
\toprule
parametro & $g_\text{min}$ Opzione A & $b$ al minimo (A) & $g_\text{min}$ Opzione B & $b$ al minimo (B) \\
\midrule
$0.009$ & $1.7\times10^{-8}$ & $2.3999$ & $0.0149$ & $2.4000$ \\
$0.073$ & $2.7\times10^{-8}$ & $2.3933$ & $0.1240$ & $2.4017$ \\
$0.148$ & $1.6\times10^{-8}$ & $2.3725$ & $0.2511$ & $2.4068$ \\
$0.300$ & $2.6\times10^{-8}$ & $2.2866$ & $0.5072$ & $2.4275$ \\
\bottomrule
\end{tabular}
\end{center}

Il quadro cambia completamente. \textbf{Opzione A}: il gap vero resta a
precisione numerica zero ($\sim10^{-8}$, compatibile con l'errore
di arrotondamento) per \emph{ogni} intensità del DM testata — il punto di incrocio si
sposta (da $b=2.4$ fino a $b\approx2.29$ per $D'=0.3$), ma resta un incrocio
\emph{vero}. Il valore positivo trovato prima a $b=2.4$ fisso era un artefatto:
misurava la distanza fra due curve nel punto sbagliato, non l'apertura reale.
\textbf{Opzione B}: il gap vero è positivo e cresce linearmente, il punto critico si
sposta di poco ($2.4\to2.43$) — qui la conclusione precedente resta valida.

\begin{center}
\includegraphics[width=0.95\textwidth]{confronto_dm_CORRETTO.png}
\end{center}
\begin{center}
\small\textit{In alto a sinistra: gap vero (minimizzato su $b$) in funzione del
parametro DM — la curva verde (Opzione A) resta a zero, la rossa (Opzione B) cresce
linearmente. In basso a destra: zoom sul punto critico di ciascuna opzione — le curve
verdi si toccano, le rosse si respingono visibilmente.}
\end{center}

\section{Perché succede: la regola di non-incrocio}

\subsection{Il risultato generale (von Neumann--Wigner)}

Per una famiglia di matrici hermitiane $H(b)$ dipendenti da un solo parametro reale
$b$, due autovalori generici \emph{tendono a evitarsi} (anticrossing) al variare di
$b$: farli coincidere esattamente richiederebbe di annullare simultaneamente più
condizioni indipendenti (per matrici reali simmetriche, tipicamente due condizioni),
un evento non generico con un solo parametro libero. Un incrocio vero (senza
repulsione) può quindi accadere \emph{solo} se una simmetria residua — un operatore
$\hat Q$ che commuta con $H(b)$ per ogni $b$ — etichetta i due stati con autovalori
diversi di $\hat Q$: in tal caso l'elemento di matrice $\langle\psi_1|H|\psi_2\rangle$
che causerebbe la repulsione è vincolato a essere \emph{esattamente} zero per ogni
$b$ (non solo approssimativamente), perché connetterebbe due settori diversi di
$\hat Q$ — proibito se $[\hat Q,H]=0$.

\subsection{Applicazione al nostro caso}

L'incrocio a $b_c$ è fra uno stato del blocco $C$ ($S_{12}=0$) e uno del blocco $A$
($S_{12}=1$) — \textbf{settori diversi} dell'operatore $\mathbf S_{12}^2$.

\begin{itemize}[leftmargin=1.4em]
\item \textbf{Opzione A}: $[\mathbf S_{12}^2, H]=0$ esattamente (verificato). L'elemento
di matrice fra lo stato di $C$ e quello di $A$ è quindi vincolato a zero per
costruzione, qualunque sia $D'$: il DM può spostare il punto in cui i due livelli si
toccano (perché altera le energie dei due rami separatamente), ma \textbf{non può mai
aprire quell'incrocio}, perché la simmetria che protegge la loro indipendenza resta
intatta.
\item \textbf{Opzione B}: $[\mathbf S_{12}^2,H]\neq0$. Non c'è più alcun numero
quantico conservato che distingua i due stati: l'elemento di matrice che li connette
non è più vincolato a zero, e generically si apre un gap — esattamente quello che si
osserva.
\end{itemize}

\textbf{Nota a margine, verificata numericamente.} Anche con l'Opzione A, DENTRO il
settore $S_{12}=1$ (dove $A$ e $B$ si mescolano, perché lì $\mathbf S^2_\text{tot}$
\emph{non} è conservato) si trova un piccolo gap reale altrove
($\sim1.6\times10^{-4}$ a $b\approx1.07$ e $D'=0.15$) — coerente con la regola:
l'Opzione A protegge \emph{solo} gli incroci fra settori diversi di $S_{12}$, non
tutti gli incroci del sistema.

\section{Conclusioni e implicazioni per la tesi}

\begin{enumerate}[leftmargin=1.4em]
\item Solo l'Opzione B (DM su tutti e tre i legami, stesso segno, proporzionale a $J$)
apre realmente il gap all'incrocio $C\to A$ che interessa il punto di lavoro scelto.
L'Opzione A (simmetrica, coerente con la teoria in forma chiusa) non lo fa mai, per
una ragione di simmetria esatta, non per una coincidenza numerica.
\item Il prezzo dell'Opzione B è la perdita completa della struttura a blocchi:
serve la diagonalizzazione $8\times8$ piena, non c'è più alcuna riduzione algebrica
disponibile.
\item Il controllo del gap va sempre fatto cercando il minimo assoluto su tutto
l'intervallo di campo, mai a un valore di campo fissato a priori — il punto critico
può spostarsi, e un valore di gap positivo calcolato nel punto sbagliato è privo di
significato fisico.
\item La spiegazione ultima è la regola di non-incrocio di von Neumann--Wigner
applicata a $\mathbf S_{12}^2$: un termine che preserva un numero quantico non può mai
aprire un incrocio fra stati che quel numero quantico distingue, indipendentemente
dalla sua intensità.
\end{enumerate}

\end{document}
